\begin{frame}{Problem Analysis}

\begin{tcolorbox}[title={문제의 목표}]
  계산 그래프가 주어질 때, 하드웨어에서 빠르게 실행하는 전략을 찾아야 합니다.
\end{tcolorbox}
\pause

\begin{itemize}
  \item 문제에서 가정하는 하드웨어는 일반적인 NPU와 비슷한 구조를 갖고 있습니다.
  \begin{itemize}
    \item 빠르고 공간이 제한된 메모리와 느리고 공간이 큰 두 종류의 메모리가 있습니다.
    \item 계산 그래프를 노드 그룹 단위로, 텐서 연산을 타일 단위로 실행합니다.
    \item 텐서 연산을 위해서는 입력 텐서의 타일을 빠른 메모리로 로드해야 합니다.
    \item 메모리 접근과 계산을 동시에 수행할 수 있습니다.
  \end{itemize}
  \pause
  \item 입력으로 어떤 계산 그래프가 주어질지 알 수 없습니다.
  \begin{itemize}
    \item 특정한 경우에 대해서만 실행 전략을 찾는 것으로는 부족합니다.
    \item 가능한 입력에 대해서 일반적으로 빠른 실행 전략을 찾는 알고리즘을 설계해야 합니다.
  \end{itemize}
\end{itemize}
\end{frame}

\begin{frame}{Problem Analysis}
\framesubtitle{Execution Strategy}

알고리즘은 크게 세 가지 변수에 대한 최적값을 찾아야 합니다.
\begin{enumerate}
  \item 하드웨어에서 한번에 실행할 노드 그룹
  \item 각 그룹에 있는 텐서 연산의 실행 타일 크기
  \item 그룹의 출력 중, 느린 메모리를 거치지 않을 텐서
\end{enumerate}
\end{frame}

\newcommand{\animframe}[1]{\resizebox{\linewidth}{!}{\includestandalone{example3-animation/#1}}}

\begin{frame}{Problem Analysis}
\framesubtitle{Example: Strategy on Residual Connection}

% ----- Colours / styles / drawing macros ONLY (no \usepackage).
% Safe to \input inside a group in the document body, so the names stay scoped.
% ======================================================================
%  tikzcommon.tex  --  Example 3 / Strategy B, software-pipelined view.
%  Each tile interval overlaps load(next) and store(prev) on the single
%  memory channel with compute, so interval = max(load+store, compute).
%  A Gantt strip at the bottom shows the staggering.  Every common element
%  is defined here once -> byte-identical across frames (no \only<> jitter).
% ======================================================================

\newcommand{\CanvasW}{16}
\newcommand{\CanvasH}{9}
\newcommand{\CanvasBB}{\useasboundingbox (-0.35,2.20) rectangle (12.7,7.7);}

% ---- palette ----------------------------------------------------------
\definecolor{ink}{RGB}{34,34,40}
\definecolor{sgA}{RGB}{37,99,235}     % subgraph 0 / compute A,B (blue)
\definecolor{sgB}{RGB}{234,88,12}     % subgraph 1 / compute C,D (orange)
\definecolor{resG}{RGB}{22,163,74}    % resident / retained (green)
\definecolor{inB}{RGB}{59,130,246}    % load  (blue)
\definecolor{outO}{RGB}{217,70,12}    % store (orange)
\definecolor{play}{RGB}{222,30,30}    % playhead cursor (red)
\definecolor{panel}{RGB}{205,208,214}
\definecolor{soft}{RGB}{246,247,249}
\definecolor{ttl}{RGB}{20,20,26}

\tikzset{
  opnode/.style ={circle,draw=ink,line width=0.9pt,fill=soft,
                  minimum size=0.82cm,inner sep=0pt,font=\small\bfseries},
  ionode/.style ={rectangle,draw=ink,line width=0.9pt,
                  fill=soft,minimum size=0.76cm,inner sep=1pt,font=\small\bfseries},
  tedge/.style  ={-{Latex[length=2.0mm]},draw=ink!75,line width=0.9pt},
  tlab/.style   ={font=\scriptsize,text=ink!65,inner sep=1pt},
  outframe/.style={draw=panel,line width=1pt},
}

% ---- output-map cells (map (7.4,5.4)-(9.1,7.1), split y=6.25) ---------
\newcommand{\tilefill}[2]{%
  \ifthenelse{\equal{#1}{top}}{\fill[#2!35] (7.4,6.25) rectangle (9.1,7.1);}{}%
  \ifthenelse{\equal{#1}{bot}}{\fill[#2!35] (7.4,5.4)  rectangle (9.1,6.25);}{}%
  \ifthenelse{\equal{#1}{both}}{\fill[#2!28] (7.4,5.4) rectangle (9.1,7.1);}{}%
}
\newcommand{\activeborder}[2]{%
  \ifthenelse{\equal{#1}{top}}%
    {\draw[#2,dash pattern=on 2.5pt off 2pt,line width=1.2pt] (7.46,6.31) rectangle (9.04,7.04);}{}%
  \ifthenelse{\equal{#1}{bot}}%
    {\draw[#2,dash pattern=on 2.5pt off 2pt,line width=1.2pt] (7.46,5.46) rectangle (9.04,6.19);}{}%
}

% ---- subgraph halos ---------------------------------------------------
% op1 (2.3,5.3) -- op2 (3.7,6.7): 45-deg axis, midpoint (3.0,6.0).
% Rectangle is centred on the midpoint then rotated about it, so it hugs
% exactly the two nodes (half-length 1.58 along axis, half-width 0.59).
\newcommand{\haloSGzero}[2]{%
  \begin{scope}[opacity=#1]
    \begin{scope}[rotate around={45:(3.0,6.0)}]
      \fill[sgA!30] (1.42,5.41) rectangle (4.58,6.59);
    \end{scope}
  \end{scope}
  \ifnum#2=1\node[font=\scriptsize\bfseries,text=sgA] at (1.95,6.98){Subgraph 0};\fi}
\newcommand{\haloSGone}[2]{%
  \begin{scope}[opacity=#1]\fill[sgB!35] (1.78,4.74) rectangle (5.30,5.86);\end{scope}
  \ifnum#2=1\node[font=\scriptsize\bfseries,text=sgB] at (4.55,4.55){Subgraph 1};\fi}

\newcommand{\sinkring}[3]{\draw[line width=1.7pt,#3] (#1,#2) circle (0.52);}
\newcommand{\recomputeTag}{%
  \node[font=\scriptsize,text=sgB] at (2.3,4.40){$\circlearrowleft$ recompute op1};
  \draw[sgB,-{Latex[length=1.6mm]},line width=0.8pt] (2.3,4.57)--(2.3,4.80);}

\newcommand{\tilemark}[1]{% #1 = formatted tensor label, e.g. $T_2$ or out
  \node[font=\scriptsize,text=ink] at (8.25,6.675){#1:t};
  \node[font=\scriptsize,text=ink] at (8.25,5.825){#1:b};}
\newcommand{\memtext}[2]{%
  \node[anchor=north west,font=\scriptsize,text=ink] at (7.4,5.18){\textbf{Fast memory:}\; #1};
  \node[anchor=north west,font=\scriptsize,text=ink] at (7.4,4.80){\textbf{Slow memory:}\; #2};}

% ======================================================================
%  Pipeline Gantt strip.  \pipe{active}  active: -1 init, 0..5 cols, 6 done
%  columns: 0 fill(ld A) | 1 K_A(ld B) | 2 K_B(ld C) | 3 K_C(ld D)
%           | 4 K_D(st C) | 5 drain(st D)
% ======================================================================
\newcommand{\pipe}[1]{% #1 = playhead x (time). <1.5 hides bar/fades all; >12.3 all solid.
  \node[font=\scriptsize\bfseries,text=ink!80,anchor=east] at (1.40,3.70){Compute};
  \node[font=\scriptsize\bfseries,text=ink!80,anchor=east] at (1.40,2.85){Memory};
  \foreach \gx in {3.3,5.1,6.9,8.7,10.5}{%
    \draw[black!16,line width=0.4pt,dash pattern=on 1.6pt off 1.6pt] (\gx,2.51)--(\gx,4.04);}
  % MEMORY channel (packed, bottleneck)
  \foreach \xL/\xR/\ml/\mc in {%
      1.5/3.3/{ld $A$:t}/inB, 3.3/5.1/{ld $A$:b}/inB,
      5.1/6.9/{ld $A$:t}/inB, 6.9/8.7/{ld $A$:b}/inB,
      8.7/10.5/{st out:t}/outO, 10.5/12.3/{st out:b}/outO}{%
    \pgfmathparse{\xL>#1 ? 0.18 : 1.0}\edef\om{\pgfmathresult}%
    \begin{scope}[opacity=\om]
      \fill[\mc!18] (\xL,2.55) rectangle (\xR,3.15);
      \draw[\mc,line width=0.7pt] (\xL,2.55) rectangle (\xR,3.15);
      \node[font=\tiny,text=ink] at ({(\xL+\xR)/2},2.85){\ml};
    \end{scope}}%
  % COMPUTE op steps: each tile after its load lands
  \foreach \xL/\opa/\opb/\cc in {%
      3.3/op1/op2/sgA, 5.1/op1/op2/sgA, 6.9/op1/op3/sgB, 8.7/op1/op3/sgB}{%
    \pgfmathparse{\xL>#1 ? 0.18 : 1.0}\edef\oA{\pgfmathresult}%
    \begin{scope}[opacity=\oA]
      \fill[\cc!22] (\xL,3.40) rectangle ({\xL+0.439},4.00);
      \draw[\cc,line width=0.7pt] (\xL,3.40) rectangle ({\xL+0.439},4.00);
      \node[font=\tiny,text=ink] at ({\xL+0.22},3.70){\opa};
    \end{scope}
    \pgfmathparse{(\xL+0.439)>#1 ? 0.18 : 1.0}\edef\oB{\pgfmathresult}%
    \begin{scope}[opacity=\oB]
      \fill[\cc!22] ({\xL+0.439},3.40) rectangle ({\xL+0.878},4.00);
      \draw[\cc,line width=0.7pt] ({\xL+0.439},3.40) rectangle ({\xL+0.878},4.00);
      \node[font=\tiny,text=ink] at ({\xL+0.659},3.70){\opb};
    \end{scope}}%
  % PLAYHEAD: vertical cursor at current time
  \pgfmathparse{(#1>=1.5)&&(#1<=12.3) ? 1 : 0}\edef\showbar{\pgfmathresult}%
  \ifdim\showbar pt>0.5pt
    \draw[play,line width=0.6pt,dash pattern=on 2.4pt off 1.8pt] (#1,2.40)--(#1,4.15);
  \fi
}

% ======================================================================
\newcommand{\StaticScene}{%
  % graph
  \node[ionode] (A)   at (1.0,5.3){$A$};
  \node[opnode] (op1) at (2.3,5.3){op1};
  \node[opnode] (op2) at (3.7,6.7){op2};
  \node[opnode] (op3) at (4.7,5.3){op3};
  \node[ionode] (out) at (6.0,5.3){out};
  \draw[tedge] (A)--(op1); \draw[tedge] (op1)--(op2);
  \draw[tedge] (op1)--(op3); \draw[tedge] (op2)--(op3); \draw[tedge] (op3)--(out);
  \node[tlab] at (2.78,6.05){$T_1$};
  \node[tlab] at (3.50,5.08){$T_1$};
  \node[tlab] at (4.42,6.05){$T_2$};
  % output map
  \draw[outframe] (7.4,5.4) rectangle (9.1,7.1);
  \draw[panel,line width=1pt] (7.4,6.25)--(9.1,6.25);
}

\begin{figure}
\centering
  \only<1>{\resizebox{\linewidth}{!}{\animframe{frame0}}}%
  \only<2>{\resizebox{\linewidth}{!}{\animframe{frame1}}}%
  \only<3>{\resizebox{\linewidth}{!}{\animframe{frame2}}}%
  \only<4>{\resizebox{\linewidth}{!}{\animframe{frame3}}}%
  \only<5>{\resizebox{\linewidth}{!}{\animframe{frame4}}}%
  \only<6>{\resizebox{\linewidth}{!}{\animframe{frame5}}}%
  \only<7>{\resizebox{\linewidth}{!}{\animframe{frame6}}}%
  \only<8>{\resizebox{\linewidth}{!}{\animframe{frame7}}}%
  \only<9>{\resizebox{\linewidth}{!}{\animframe{frame8}}}%
  \only<10>{\resizebox{\linewidth}{!}{\animframe{frame9}}}%
  \only<11>{\resizebox{\linewidth}{!}{\animframe{frame10}}}%
  \only<12>{\resizebox{\linewidth}{!}{\animframe{frame11}}}
\end{figure}
\end{frame}

\begin{frame}{Problem Analysis}
\framesubtitle{이거 웰노운 아닌가?}

\begin{itemize}
  \item 결국 문제의 핵심은 operator fusion을 잘하는 것입니다.
  \begin{itemize}
    \item Operator fusion은 머신러닝 시스템 분야에서 활발히 연구되고 있는 주제입니다.
    \item PyTorch, JAX, TensorRT 등 다양한 런타임에 구현되어 있습니다.
    \item ``코딩할 시간도 얼마 없는데 이미 잘 구현된 아이디어를 참고하면 되지 않을까?''
  \end{itemize}
  \pause
  \item 대부분의 런타임에서 operator fusion은 rule based 기법으로 구현합니다.
  \begin{itemize}
    \item 자주 나타나는 패턴과 최적화 방식을 미리 입력해 두고, 패턴 매칭을 통해 최적화합니다.
    \item 예시: 행렬 곱의 결과를 ReLU가 사용하는 경우, 두 연산을 함께 실행합니다.
  \end{itemize}
\end{itemize}

\pause
\begin{figure}
  \centering
  \only<1-3>{\resizebox{!}{2.5cm}{\includegraphviz[width=\linewidth]{graphs/nn-twolayer}}}%
  \only<4>{\resizebox{!}{2.5cm}{\includegraphviz[width=\linewidth]{graphs/nn-twolayer-fused}}}
\end{figure}
\end{frame}

\begin{frame}{Problem Analysis}
\framesubtitle{날로 먹기란 쉽지 않다}

\begin{itemize}
  \item Rule based 기법을 적용하려면 패턴을 그래프에서 찾기 쉬워야 합니다.
  \begin{itemize}
    \item 일반적인 그래프는 행렬 곱, ReLU, reshape 등 다양한 종류의 연산이 있습니다.
    \pause
    \item 하지만 이 문제의 그래프에서는 연산 종류가 matmul, elementwise 두 종류뿐입니다.
    \item 원래는 연산 종류로 실행 시간을 예상할 수 있지만, 문제에서는 임의로 주어집니다.
  \end{itemize}
\end{itemize}

\pause
\begin{figure}
  \centering
  \includegraphviz[width=0.9\linewidth]{graphs/arbitrary-cost}
\end{figure}

\pause
\begin{tcolorbox}[title={전체적인 전략}]
  패턴 매칭 대신, 다양한 최적화 전략을 탐색하는 알고리즘을 구현해야 합니다.
\end{tcolorbox}
\end{frame}